\subsubsection*{04.05.01 Projektdesign}
\phantomsection
\addcontentsline{toc}{subsubsection}{04.05.01 Projektdesign}
\label{sec:04.05.01_Projektdesign}

%%% Situation 

\par Die Ausgangslage für das Projektdesign war durch die Tragweite der Transformation geprägt: Alle im \gls{r3} geführten \gls{bk} sollten innerhalb von acht Jahren aus dem historisch gewachsenen, heterogenen \gls{r3}-System in ein modernes, standardisiertes \gls{s4}-System überführt werden. Das Umfeld war gekennzeichnet durch zahlreiche Stakeholder, komplexe Abhängigkeiten der über 400 zu entwickelnden \gls{ricefw}-Objekte und die Forderung, den laufenden Betrieb während der Transformation aufrechtzuerhalten. Zusätzliche Scope-Anpassungen wie die Neueinführung von SAP als \gls{erp} bei \gls{adn}, die Nutzung von \gls{mbc} für Bankkommunikation und die Einführung von \gls{gls-drc} zur Sicherstellung der \gls{gls-ere} im \gls{b2b}-Umfeld verstärkten die Komplexität. In diesem Kontext war ein Projektdesign erforderlich, das technische, organisatorische, prozessuale und kulturelle Herausforderungen gleichermaßen adressierte, um die Transformation erfolgreich umzusetzen.